% !PRINTED 356
\weberpage{356}
%<<CONT
of such equations, however, we obtain two more by multiplying (1) by $x^{-m}$ and by $x^{-m-n}$. This gives
\[
(3)\quad r^{n+m} = \frac{-g\sin m\vartheta}{\sin n\vartheta},
\]
\[
(4)\quad r^m = \frac{-g\sin(n+m)\vartheta}{e\sin n\vartheta},
\]
and of these three equations each follows from the other two. If one eliminates $r$ from two of them, there results, if, as above,
\[
\lambda = \frac{g^n}{e^{n+m}}
\]
is set,
\[
(5)\quad \lambda = (-1)^n \frac{\sin n\vartheta^n \sin m\vartheta^m}{\sin(n+m)\vartheta^{n+m}}.
\]

Since one need compute only one of the two conjugate roots, one may restrict $\vartheta$ to the first quadrant, but must then in some circumstances also take $r$ negative.

How one computes from formula (5), by means of the trigonometric tables, the angle $\vartheta$, and then from one of the three formulas (2), (3), (4) the corresponding value of $r$, we shall now demonstrate on the example considered above
\[
x^3 - 2x - 2 = 0
\]
as well.

Here equation (5) becomes
\[
(6)\quad \frac{1}{2} = \frac{\sin 2\vartheta^2 \sin\vartheta}{\sin 3\vartheta^3},
\]
or
\[
(7)\quad 3\log\sin 3\vartheta - 2\log\sin 2\vartheta - \log\sin\vartheta = 0{,}3010300,
\]
from which it is first of all evident that the angle $\vartheta$ lies in the interval from 0 to $60^0$. After a few trials one finds first that $\vartheta$ lies between $33^0$ and $34^0$, and if one then computes for these two values the difference
\[
(8)\quad 3\log\sin 3\vartheta - 2\log\sin 2\vartheta - \log\sin\vartheta - 0{,}3010300
\]
one obtains first, to three decimals,
\[
0{,}026, \quad -0{,}013,
\]
and from this, by interpolation, there results the more exact value
\[
\vartheta = 33^0\,40'.
\]
